\documentclass[landscape,final]{baposter}

\usepackage{times}
\usepackage{calc}
\usepackage{graphicx}
\usepackage{amsmath}
\usepackage{amssymb}
\usepackage{relsize}
\usepackage{multirow}
\usepackage{bm}
\usepackage{psfrag}

\usepackage{graphicx}
\usepackage{multicol}

\usepackage{pgfbaselayers}
\pgfdeclarelayer{background}
\pgfdeclarelayer{foreground}
\pgfsetlayers{background,main,foreground}

\usepackage{helvet}
%\usepackage{bookman}
\usepackage{palatino}

\newcommand{\captionfont}{\footnotesize}

\selectcolormodel{cmyk}

\graphicspath{{images/}}

%%%%%%%%%%%%%%%%%%%%%%%%%%%%%%%%%%%%%%%%%%%%%%%%%%%%%%%%%%%%%%%%%%%%%%%%%%%%%%%%
% Multicol Settings
%%%%%%%%%%%%%%%%%%%%%%%%%%%%%%%%%%%%%%%%%%%%%%%%%%%%%%%%%%%%%%%%%%%%%%%%%%%%%%%%
\setlength{\columnsep}{0.7em}
\setlength{\columnseprule}{0mm}


%%%%%%%%%%%%%%%%%%%%%%%%%%%%%%%%%%%%%%%%%%%%%%%%%%%%%%%%%%%%%%%%%%%%%%%%%%%%%%%%
% Save space in lists. Use this after the opening of the list
%%%%%%%%%%%%%%%%%%%%%%%%%%%%%%%%%%%%%%%%%%%%%%%%%%%%%%%%%%%%%%%%%%%%%%%%%%%%%%%%
\newcommand{\compresslist}{%
\setlength{\itemsep}{1pt}%
\setlength{\parskip}{0pt}%
\setlength{\parsep}{0pt}%
}


%%%%%%%%%%%%%%%%%%%%%%%%%%%%%%%%%%%%%%%%%%%%%%%%%%%%%%%%%%%%%%%%%%%%%%%%%%%%%%
%%% Begin of Document
%%%%%%%%%%%%%%%%%%%%%%%%%%%%%%%%%%%%%%%%%%%%%%%%%%%%%%%%%%%%%%%%%%%%%%%%%%%%%%

\begin{document}

%%%%%%%%%%%%%%%%%%%%%%%%%%%%%%%%%%%%%%%%%%%%%%%%%%%%%%%%%%%%%%%%%%%%%%%%%%%%%%
%%% Here starts the poster
%%%---------------------------------------------------------------------------
%%% Format it to your taste with the options
%%%%%%%%%%%%%%%%%%%%%%%%%%%%%%%%%%%%%%%%%%%%%%%%%%%%%%%%%%%%%%%%%%%%%%%%%%%%%%
\typeout{Poster Starts}
\definecolor{blue}{cmyk}{0.9,0,0,0.0}
\definecolor{reddishblue}{cmyk}{1,0.22,0,0.0}
\definecolor{black}{cmyk}{0,0,0.0,1.0}

\definecolor{lightblue}{cmyk}{.3,0,0,0.0}
\definecolor{lighterblue}{cmyk}{.2,0,0,0.0}
\definecolor{lightestblue}{cmyk}{.1,0,0,0.0}
\begin{poster}{
  % Show grid to help with alignment
  grid=no,
  % Column spacing
  colspacing=1em,
  % Color style
  bgColorOne=lighterblue,
  bgColorTwo=lightestblue,
  borderColor=reddishblue,
  headerColorOne=blue,
  headerColorTwo=reddishblue,
  headerFontColor=black,
  boxColorOne=lightblue,
  boxColorTwo=lighterblue,
  % Format of textbox
  textborder=roundedleft,
  % Format of text header
  eyecatcher=yes,
  headerborder=open,
  headerheight=0.08\textheight,
  headershape=roundedright,
  headershade=plain,
  headerfont=\Large\textsf, %Sans Serif
  boxshade=plain,
  background=shade-tb,
  %background=plain,
  linewidth=2pt
  }
  % Eye Catcher
  {{\begin{minipage}{7em}
    \includegraphics[height=5em]{amllogo}
  \end{minipage}}
  }
  % Title
  {\sf %Sans Serif
  %\bf% Serif
  Automatic Topic Discovery in 100 Years of General Conference Talks}
  % Authors
  {\sf %Sans Serif
  % Serif
  Matthew Gardner
  (mjg82@byu.edu)
  and
  Eric Ringger
  (ringger@cs.byu.edu)
  }
  % University logo
  {{\begin{minipage}{7em}
    \hfill
    \includegraphics[height=5em]{byunlp}
  \end{minipage}}
  }

  \tikzstyle{light shaded}=[top color=baposterBGtwo!30!white,bottom color=baposterBGone!30!white,shading=axis,shading angle=30]

  % Width of left inset image
     \newlength{\leftimgwidth}
     \setlength{\leftimgwidth}{0.78em+8.0em}

%%%%%%%%%%%%%%%%%%%%%%%%%%%%%%%%%%%%%%%%%%%%%%%%%%%%%%%%%%%%%%%%%%%%%%%%%%%%%%
%%% Now define the boxes that make up the poster
%%%---------------------------------------------------------------------------
%%% Each box has a name and can be placed absolutely or relatively.
%%% The only inconvenience is that you can only specify a relative position 
%%% towards an already declared box. So if you have a box attached to the 
%%% bottom, one to the top and a third one which should be in between, you 
%%% have to specify the top and bottom boxes before you specify the middle 
%%% box.
%%%%%%%%%%%%%%%%%%%%%%%%%%%%%%%%%%%%%%%%%%%%%%%%%%%%%%%%%%%%%%%%%%%%%%%%%%%%%%
    %
    % A coloured circle useful as a bullet with an adjustably strong filling
    \newcommand{\colouredcircle}[1]{%
      \tikz{\useasboundingbox (-0.2em,-0.32em) rectangle(0.2em,0.32em); \draw[draw=black,fill=baposterBGone!80!black!#1!white,line width=0.03em] (0,0) circle(0.18em);}}

%%%%%%%%%%%%%%%%%%%%%%%%%%%%%%%%%%%%%%%%%%%%%%%%%%%%%%%%%%%%%%%%%%%%%%%%%%%%%%
  \headerbox{Contribution}{name=contribution,column=0,row=0}{
%%%%%%%%%%%%%%%%%%%%%%%%%%%%%%%%%%%%%%%%%%%%%%%%%%%%%%%%%%%%%%%%%%%%%%%%%%%%%%
   {}
   
   We use state of the art natural language processing techniques to gain new
   insights from 100 years of General Conference Talks.  Using computers to
   sift through a corpus that would be overwhelmingly large to process
   manually, we can discover relationships between talks and trends over time
   that help us to better appreciate the trove of knowledge in these inspired
   discourses.

 }

%%%%%%%%%%%%%%%%%%%%%%%%%%%%%%%%%%%%%%%%%%%%%%%%%%%%%%%%%%%%%%%%%%%%%%%%%%%%%%
  \headerbox{Topic Discovery with LDA}
  {name=lda,column=0,below=contribution}{

	  We automatically discover topics using Latent Dirichlet Allocation (LDA).
	  The LDA algorithm works as follows:
	  \begin{itemize}
		\item There is some number of topics in the corpus (we picked 100).
		\item Assume each word in each document comes from one topic.
		\item Use statistical relationships among words in the corpus to group 
		  the words into topics.
	  \end{itemize}

  }
%%%%%%%%%%%%%%%%%%%%%%%%%%%%%%%%%%%%%%%%%%%%%%%%%%%%%%%%%%%%%%%%%%%%%%%%%%%%%%
%%%%%%%%%%%%%%%%%%%%%%%%%%%%%%%%%%%%%%%%%%%%%%%%%%%%%%%%%%%%%%%%%%%%%%%%%%%%%%
  \headerbox{LDA Graphical Model}
  {name=methodology,column=0,below=lda}{
%%%%%%%%%%%%%%%%%%%%%%%%%%%%%%%%%%%%%%%%%%%%%%%%%%%%%%%%%%%%%%%%%%%%%%%%%%%%%%

	\begin{centering}
	  \psfrag{b}[C][C]{\Large$\beta$}
	  \psfrag{a}[C][C]{\Large$\alpha$}
	  \psfrag{t}[C][C]{\Large$\theta$}
	  \psfrag{f}[C][C]{\Large$\phi$}
	  \psfrag{w}[C][C]{\Large$w$}
	  \psfrag{z}[C][C]{\Large$z$}
	  \psfrag{N}[C][C]{\large$N$}
	  \psfrag{M}[C][C]{\large$M$}
	  \psfrag{K}[C][C]{\large$K$}
	  \includegraphics[height=.485\linewidth]{lda_model}
	  \begin{align*} 
		M &= \mathrm{Number\ of\ documents} \\
		N &= \mathrm{Number\ of\ words\ in\ a\ document} \\
		K &= \mathrm{Number\ of\ topics} \\
		w &= \mathrm{Word\ in\ a\ document} \\
		z &= \mathrm{Topic\ corresponding\ to\ word\ } w \\
		w_i|z_i,\phi^{(z_i)} &\sim \mathrm{Categorical}(\phi^{(z_i)}) \\
		\phi &\sim \mathrm{Dirichlet}(\beta) \\
		z_i|\theta^{(d_i)} &\sim \mathrm{Categorical}(\theta^{(d_i)}) \\ 
		\theta &\sim \mathrm{Dirichlet}(\alpha) 
	  \end{align*}
	\end{centering}

  }

%%%%%%%%%%%%%%%%%%%%%%%%%%%%%%%%%%%%%%%%%%%%%%%%%%%%%%%%%%%%%%%%%%%%%%%%%%%%%%
  \headerbox{Results: Topic--Word Distributions}
  {name=words,column=1,span=2,row=0}{
%%%%%%%%%%%%%%%%%%%%%%%%%%%%%%%%%%%%%%%%%%%%%%%%%%%%%%%%%%%%%%%%%%%%%%%%%%%%%%

	The main result of using LDA is a list of which words are most likely to be
	used in each topic.  In the algorithm, the topics just have numbers; the
	computer merely looks for statistical relationships between words.  Even
	so, the topics are easily distinguishable.  We present here a few of the
	topics we found.

	\begin{center}
	  \begin{tabular}{cccccccc}
		Topic 19&Topic 33&Topic 52&Topic 37&Topic 7&Topic 26&Topic 28&Topic 22\\
		\hline
		tithing&welfare&women&children&government&see&years&true \\
		pay&poor&society&home&united&d&one&honest \\
		money&church&relief&family&liberty&c&two&good \\
		debt&fast&sisters&parents&nation&1&hundred&integrity \\
		means&help&hinckley&teach&states&2&five&virtue \\
		out&services&each&child&country&3&three&character \\
		people&program&b&families&freedom&s&thousand&honesty \\
		one&plan&gordon&homes&land&ne&twenty&without \\
		paid&care&lives&love&free&4&ago&one \\
		poor&needs&service&mother&men&5&four&men \\
		tithes&food&sister&taught&constitution&alma&year&honor \\
		get&need&forward&fathers&people&6&six&respect \\
		more&family&strength&responsibility&america&7&over&god \\
	  \end{tabular}
	\end{center}

  }

%%%%%%%%%%%%%%%%%%%%%%%%%%%%%%%%%%%%%%%%%%%%%%%%%%%%%%%%%%%%%%%%%%%%%%%%%%%%%%
  \headerbox{Results: Topics over Time}
  {name=results topics over time,column=1,span=2,below=words}{
%%%%%%%%%%%%%%%%%%%%%%%%%%%%%%%%%%%%%%%%%%%%%%%%%%%%%%%%%%%%%%%%%%%%%%%%%%%%%%

	Because each talk has a date associated with it, we can use LDA to analyze
	the distribution of these topics over time, seeing at what points certain
	topics were most discussed.  In doing this, we see clear correlations with
	real-world occurrences.

	\begin{center}
	  \includegraphics[height=.305\linewidth]{topic19}
	  \includegraphics[height=.305\linewidth]{topic33}

	  \includegraphics[height=.305\linewidth]{topic52}
	  \includegraphics[height=.305\linewidth]{topic56}
	\end{center}

  }

%%%%%%%%%%%%%%%%%%%%%%%%%%%%%%%%%%%%%%%%%%%%%%%%%%%%%%%%%%%%%%%%%%%%%%%%%%%%%%
  \headerbox{Results: Topics by Speaker}
  {name=speaker,column=3,span=1}{
%%%%%%%%%%%%%%%%%%%%%%%%%%%%%%%%%%%%%%%%%%%%%%%%%%%%%%%%%%%%%%%%%%%%%%%%%%%%%%

	Because each talk has a speaker, we can look at topics to see who was most
	likely to use them in their talks.  Again we see here that our automatic
	analysis is validated by real-world occurrences.

	\begin{center}
	  \begin{tabular}{cc}
		Relief Society Topic&Sunday School Topic\\
		\hline
		Sister Barbara B. Smith&Karl G. Maeser \\
		Mary Ellen W. Smoot&George Reynolds \\
		Mary Ellen Smoot&E. C. Phillips \\
		Elaine L. Jack&Josiah H. Burrows \\
		Marian R. Boyer&George Goddard \\
		Virginia U. Jensen&Joseph J. Jackson\\
		Wayne M. Hancock&L. John Nuttall \\
		Shirley W. Thomas&Sunday School Union \\
		Bonnie D. Parkin&John D. Peters \\
		Kathleen H. Hughes&Joseph Jackson \\
		Sheri L. Dew&George Q. Cannon \\
		Richard H. Winkel&Heber J. Grant \\
		Aileen H. Clyde&Joseph Summerhays \\
	  \end{tabular}
	\end{center}

  }

%%%%%%%%%%%%%%%%%%%%%%%%%%%%%%%%%%%%%%%%%%%%%%%%%%%%%%%%%%%%%%%%%%%%%%%%%%%%%%
  \headerbox{Future Work}{name=future work,column=3,span=1,below=speaker}{
%%%%%%%%%%%%%%%%%%%%%%%%%%%%%%%%%%%%%%%%%%%%%%%%%%%%%%%%%%%%%%%%%%%%%%%%%%%%%%

	We have used natural language processing techniques on a novel data
	source---LDS General Conference talks.  However, our exploration of the
	results has been quite preliminary.  Future areas of research could 
	include:
	\begin{itemize}
	  \item A more thorough analysis of religious discourse in the LDS Church,
		answering questions such as ``Which topics are constant over time, and
		which come and go?''
	  \item Explore other religious and scholarly LDS texts, including the
		Journal of Discourses, the Book of Mormon, the Doctrine and Covenants,
		the collected works of Hugh Nibley, or the past 50 years of BYU
		Studies.
	  \item Not only explore other data sets in isolation, but compare them;
		see how the past 100 years of General Conference differs from the
		previous 100 years in the Journal of Discourses.
	  \item The use of other natural language processing techniques to analyze
		the data.  Possibilities include Topical N-grams, an extension of LDA
		that returns strings of words in the topic instead of isolated words.
	\end{itemize}

  }

\end{poster}

\end{document}
